% ============================================================================
%  E/21-lo-trinh-va-tai-nguyen — file chương
%  Ngày tạo: 2026-09-07 | Sửa gần nhất: 2026-09-07 | Ôn gần nhất: chưa
%  Dependency: E/19, E/20. Mức độ: cốt lõi.
% ============================================================================
\documentclass[../../english.tex]{subfiles}
\begin{document}
\chapter{Lộ trình mười hai tuần và tài nguyên (A Twelve-Week Plan)}
\ptrtarget{E/21}\ptrtarget{E/21-lo-trinh-va-tai-nguyen}
\dependency{\ptr{E/19} cho bốn nhánh; \ptr{E/20} cho bộ công cụ. Chương này gộp cả quyển thành lịch.}

\begin{tomtat}
Chương cuối biến toàn bộ quyển sách thành một lịch có thể bắt đầu vào thứ Hai tuần này. Nội dung gồm: một tuần mẫu năm giờ chia theo bốn nhánh; lộ trình mười hai tuần chia ba giai đoạn, mỗi giai đoạn có một trọng tâm và một sản phẩm cụ thể; bốn chỉ số đo được để phân biệt tiến bộ thật với cảm giác bận rộn; cách xử lý khi bạn lỡ nhịp; và danh mục tài nguyên đã chọn lọc theo từng nhánh. Nếu chỉ nhớ một điều: hãy chọn mức tối thiểu đủ thấp để bạn làm được vào ngày tệ nhất, vì cái quyết định kết quả sau một năm là số tuần bạn \emph{không} bỏ, chứ không phải cường độ của tuần tốt nhất.
\end{tomtat}

\begin{nentang}
\kn{mức tối thiểu}{lượng công việc nhỏ nhất bạn cam kết làm kể cả ngày bận nhất --- để chuỗi không đứt.}
\kn{sản phẩm giai đoạn}{một thứ cụ thể hoàn thành ở cuối mỗi bốn tuần: một bản tóm tắt, một phần bài, một bảng thuật ngữ.}
\kn{chỉ số tiến bộ}{con số đo được thay cho cảm giác: số từ viết trong 15 phút, số lần khựng trên một trang.}
\kn{lỡ nhịp}{bỏ một hoặc vài buổi --- việc chắc chắn sẽ xảy ra, nên phải có quy tắc xử lý sẵn.}
\kn{ôn theo lịch giãn}{ôn lại các chương của quyển này theo khoảng cách tăng dần thay vì đọc lại tuần tự.}
\kn{ngân hàng tài nguyên}{danh mục nguồn đã chọn lọc, xếp theo nhánh và theo mục đích.}
\end{nentang}

\begin{khoigoi}
\h{Một tuần luyện tập cân bằng trông như thế nào, nếu bạn chỉ có năm giờ?}
\h{Làm sao phân biệt tiến bộ thật với cảm giác đang bận rộn học hành?}
\h{Bạn bỏ mất một tuần vì công việc. Nên làm bù hay bỏ qua?}
\h{Sau mười hai tuần thì làm gì tiếp?}
\h{Nên ôn lại quyển sách này thế nào cho hiệu quả?}
\end{khoigoi}

Hai chương trước cho bạn khung và công cụ. Chương này trả lời câu hỏi còn lại: làm gì vào thứ Hai tuần này.

Một lưu ý trước khi vào lịch. Lộ trình dưới đây được thiết kế cho \emph{năm giờ một tuần} --- không phải vì đó là con số tối ưu mà vì đó là con số mà một người đi làm duy trì được trong nhiều tháng. Một lộ trình mười lăm giờ một tuần trông ấn tượng hơn trên giấy và bị bỏ sau ba tuần, và ba tuần cường độ cao thua xa sáu tháng cường độ vừa.

\section{Một tuần mẫu}

Câu hỏi mở đầu thứ nhất. Năm giờ, chia theo bốn nhánh của \ptr{E/19}.

\begin{center}\footnotesize
\rowcolors{2}{white}{tblalt}
\begin{longtable}{@{}L{2.2cm}L{1.6cm}L{5.0cm}L{5.4cm}@{}}
\toprule
\textbf{Nhánh} & \textbf{Thời gian} & \textbf{Việc cụ thể} & \textbf{Mức tối thiểu ngày bận}\\
\midrule
\endfirsthead
\toprule
\textbf{Nhánh} & \textbf{Thời gian} & \textbf{Việc cụ thể} & \textbf{Mức tối thiểu ngày bận (tiếp)}\\
\midrule
\endhead
1. Tiếp nhận dễ & 75 phút & Đọc bài tổng quan hoặc giáo trình ở mức hiểu gần hết; xem một bài giảng & Đọc một trang, không tra từ nào\\
2. Học chủ đích & 75 phút & Xây bảng thuật ngữ (\ptr{B/10}); ghi cụm (\ptr{D/17}); tra bằng ngữ liệu & Thêm ba mục vào bảng thuật ngữ\\
3. Sản xuất & 90 phút & Nhật ký kỹ thuật 15 phút \(\times\) 4; viết lại một đoạn theo \ptr{C/15} & Viết năm câu, không dừng để tra\\
4. Lưu loát & 60 phút & Đọc nhanh có bấm giờ; nói lại một bài báo trong 3 phút & Ba phút nói lại nội dung vừa đọc\\
\bottomrule
\end{longtable}
\end{center}

Cột cuối là cột quan trọng nhất và cũng là cột hay bị bỏ khi người ta lập kế hoạch. Nó tồn tại vì một lý do đã nói ở \ptr{E/19} mục 5: thứ giết thói quen không phải làm ít mà là \emph{đứt chuỗi}. Một ngày làm năm phút giữ chuỗi nguyên vẹn; một ngày bỏ hẳn mở đường cho ngày thứ hai.

\begin{cannho}
Ghép mỗi nhánh vào một khoảng thời gian cố định trong ngày, gắn với một việc bạn vốn đã làm. Ví dụ: nhánh 3 ngay sau khi mở máy buổi sáng, trước khi đọc email; nhánh 2 trong lúc đọc bài báo cho công việc. Việc gắn vào một mốc sẵn có hiệu quả hơn nhiều so với đặt lịch nhắc, vì mốc sẵn có luôn xảy ra còn lời nhắc thì bị tắt đi.
\end{cannho}

\section{Mười hai tuần, ba giai đoạn}

Mười hai tuần chia thành ba giai đoạn bốn tuần, mỗi giai đoạn có một trọng tâm và kết thúc bằng một sản phẩm cụ thể. Sản phẩm quan trọng: nó biến ``học tiếng Anh'' thành một việc có điểm kết thúc quan sát được.

\begin{figure}[H]\centering
\begin{tikzpicture}[node distance=0.45cm,
  g/.style={draw=steel,fill=bluebg,rounded corners=3pt,inner sep=6pt,align=left,font=\small,text width=10.5cm},
  ar/.style={-{Latex[length=2.4mm]},draw=accent,thick},
  lb/.style={font=\scriptsize\itshape,color=tiercol,align=left,text width=2.4cm}]
\node[g] (g1) {\textbf{Tuần 1--4 --- Nền và đọc.} Trọng tâm: \ptr{A/02}, \ptr{B/05}, \ptr{B/06}, \ptr{B/10}.\\ \emph{Sản phẩm:} bảng thuật ngữ 60 mục \(+\) một bản đồ lập luận của một bài báo.};
\node[g,below=of g1] (g2) {\textbf{Tuần 5--8 --- Viết rõ.} Trọng tâm: \ptr{C/11}, \ptr{C/12}, \ptr{C/13}, \ptr{C/14}.\\ \emph{Sản phẩm:} một bản tóm tắt 200 từ \(+\) một phần mở bài theo ba nước đi CARS.};
\node[g,below=of g2] (g3) {\textbf{Tuần 9--12 --- Sắc thái và tự động hoá.} Trọng tâm: \ptr{D/16}, \ptr{D/17}, \ptr{C/15}.\\ \emph{Sản phẩm:} ngân hàng 40 cụm \emph{đã dùng} \(+\) một bài viết hoàn chỉnh đã sửa bốn lượt.};
\draw[ar] (g1) -- (g2); \draw[ar] (g2) -- (g3);
\node[lb,right=0.2cm of g1.east] {đọc trước, vì viết cần nguyên liệu};
\node[lb,right=0.2cm of g3.east] {mỗi cụm phải \emph{được dùng}};
\end{tikzpicture}
\caption{Ba giai đoạn. Thứ tự không tuỳ tiện: đọc đi trước vì viết cần nguyên liệu và cần mẫu; sắc thái đi sau vì nó chỉ có ý nghĩa khi bạn đã có bài để áp vào.}
\label{fig:lo-trinh-12-tuan}
\end{figure}

Trong cả mười hai tuần, bốn nhánh vẫn chạy song song theo tuần mẫu ở mục 1. Cái thay đổi giữa các giai đoạn là \emph{nội dung} của nhánh 2 và nhánh 3, chứ không phải cấu trúc tuần.

\section{Bốn chỉ số}

Câu hỏi mở đầu thứ hai là câu hỏi quan trọng nhất của chương, vì cảm giác tiến bộ trong ngôn ngữ rất không đáng tin theo cả hai chiều: bạn có thể cảm thấy trì trệ trong lúc đang tiến, và cảm thấy hiệu quả trong lúc chỉ đang bận rộn.

Bốn chỉ số dưới đây đo được, mất chưa tới mười phút mỗi tuần, và mỗi chỉ số ứng với một nhánh.

\begin{center}\footnotesize
\rowcolors{2}{white}{tblalt}
\begin{longtable}{@{}L{4.4cm}L{1.3cm}L{4.4cm}L{4.1cm}@{}}
\toprule
\textbf{Chỉ số} & \textbf{Nhánh} & \textbf{Cách đo} & \textbf{Hướng kỳ vọng}\\
\midrule
\endfirsthead
\toprule
\textbf{Chỉ số} & \textbf{Nhánh} & \textbf{Cách đo} & \textbf{Hướng kỳ vọng (tiếp)}\\
\midrule
\endhead
Số lần khựng trên một trang & 1 & Đọc một trang tài liệu ngành, đếm chỗ phải dừng & Giảm dần; dưới hai là ngưỡng trôi chảy\\
Số mục mới trong bảng thuật ngữ & 2 & Đếm cuối tuần & Cao lúc đầu, giảm dần khi phủ hết ngành\\
Số từ viết trong 15 phút không dừng & 3 & Bấm giờ, đếm từ & Tăng; đây là chỉ số nhạy nhất\\
Số cụm \emph{đã dùng} trong bài thật & 3 \(+\) 4 & Đếm trong ngân hàng cụm & Tăng; khác hẳn số cụm đã ghi\\
\bottomrule
\end{longtable}
\end{center}

Chỉ số thứ ba đáng chú ý vì nó phản ứng nhanh nhất. Nếu bạn đang bỏ trống nhánh 3 và bắt đầu viết mười lăm phút mỗi ngày, con số này thường tăng rõ trong ba tới bốn tuần --- sớm hơn nhiều so với bất kỳ dấu hiệu nào khác. Đó cũng là lý do nó đáng đo: nó cho bạn bằng chứng sớm rằng cách làm đang hoạt động, đúng lúc động lực dễ tụt nhất.

Chỉ số thứ tư là chỉ số duy nhất phân biệt được học thật với sưu tầm. Số cụm \emph{đã ghi} có thể tăng vô hạn mà không đổi gì; số cụm \emph{đã dùng} thì không.

\section{Khi lỡ nhịp}

Câu hỏi mở đầu thứ ba. Bạn sẽ lỡ --- không phải nếu mà là khi. Có sẵn quy tắc thì lỡ một tuần không thành bỏ cuộc.

\emph{Quy tắc một --- không làm bù.} Đây là quy tắc quan trọng nhất và đi ngược trực giác. Làm bù biến việc lỡ nhịp thành một khoản nợ, và nợ tích lại làm người ta bỏ hẳn. Buổi đã lỡ thì mất; tiếp tục từ hôm nay.

\emph{Quy tắc hai --- hạ xuống mức tối thiểu, đừng bỏ.} Trong tuần bận, chạy cột cuối của bảng ở mục 1. Năm câu một ngày vẫn giữ được thói quen; và quan trọng hơn, nó giữ được việc \emph{tự nhận mình là người đang luyện}, thứ mà một tuần trống hoàn toàn phá mất.

\emph{Quy tắc ba --- nếu lỡ hơn hai tuần, đừng quay lại chỗ cũ.} Bắt đầu lại từ đầu giai đoạn hiện tại với mức tối thiểu trong một tuần, rồi mới lên lại cường độ đầy đủ. Quay thẳng về cường độ cũ sau một quãng nghỉ dài gần như luôn dẫn tới lần bỏ thứ hai.

\begin{cambay}
Cái bẫy tinh vi nhất không phải bỏ cuộc mà là \emph{chuyển sang nhánh dễ}. Khi mệt, người ta có xu hướng dồn hết vào nhánh 1 --- đọc và nghe --- vì nó nhẹ nhất và vẫn cho cảm giác đang học. Nếu bạn thấy tuần này mình chỉ đọc và không viết dòng nào, đó không phải một tuần luyện nhẹ; đó là một tuần mất cân bằng, và nó chính là trạng thái mà \ptr{E/19} mô tả như nguyên nhân của việc dừng lại ở một mức.
\end{cambay}

\section{Sau mười hai tuần và cách ôn lại quyển này}

Câu hỏi mở đầu thứ tư và thứ năm.

Sau mười hai tuần, cấu trúc tuần vẫn giữ nguyên nhưng \emph{nội dung} nên gắn với một mục tiêu thật: một bài báo bạn đang viết, một buổi trình bày, một chương luận án. Lý do không phải động lực mà là chất lượng phản hồi --- một bài thật có người phản biện thật, và phản hồi thật là thứ mà mọi bài luyện tự tạo đều thiếu (\ptr{E/19} mục 4).

Về việc ôn lại quyển sách: đừng đọc lại tuần tự. Cách hiệu quả hơn, dựa trên hai kết quả đã nhắc ở \ptr{E/19} --- truy xuất củng cố mạnh hơn đọc lại\cite{roediger2006} và ôn giãn cách hơn ôn dồn\cite{cepeda2006} --- gồm ba bước.

\emph{Bước một --- ôn bằng phần \emph{Ôn nhanh}, không bằng thân chương.} Che cột đáp án, tự trả lời, rồi mới mở ra. Mỗi mục có ghi mức ưu tiên; lần ôn đầu chỉ làm mức 3.

\emph{Bước hai --- ôn theo lịch giãn.} Một chương ôn sau một tuần, rồi sau ba tuần, rồi sau hai tháng. Ghi ngày ôn vào dòng ``Ôn gần nhất'' ở đầu mỗi tệp chương.

\emph{Bước ba --- ôn theo triệu chứng, không theo thứ tự.} Khi gặp một vấn đề thật --- người phản biện nói bài khó theo, bạn bí ở phần mở bài --- hãy dùng bảng tra ngược ở đầu mỗi phần để nhảy thẳng tới chương liên quan. Đây là cách ôn có giá trị nhất vì nó xảy ra đúng lúc bạn có nhu cầu.

\section{Ngân hàng tài nguyên}

Danh mục dưới đây xếp theo nhánh. Nguyên tắc chọn: mỗi mục đều dùng được ngay và không trùng chức năng với mục khác.

\emph{Nhánh 1 --- tiếp nhận dễ.} Nguồn tốt nhất là bài tổng quan và giáo trình \emph{trong chính ngành bạn}, vì chúng có mật độ từ mới thấp mà vẫn đúng miền khái niệm. Về nguyên tắc chọn tài liệu ở đúng mức, xem\cite{daybamford1998}. Với nghe, các bài giảng hội nghị có phụ đề và bản ghi là lựa chọn tự nhiên (\ptr{B/09}).

\emph{Nhánh 2 --- học chủ đích.} Từ điển cho người học\cite{oald,ldoce}; ngữ liệu\cite{coca}; từ điển gốc từ\cite{etymonline}; danh sách từ học thuật\cite{coxhead2000} và danh sách theo tần suất\cite{ngsl}. Về nguyên lý học từ vựng, quyển tham chiếu là\cite{nation2013}, bản ngắn hơn và thực hành hơn là\cite{webbnation2017}; về mặt tâm lý học của việc học từ, xem\cite{schmitt2000}. Cho ngữ pháp, dùng như từ điển tra cứu chứ không đọc tuần tự\cite{swan2016,murphy2019}.

\emph{Nhánh 3 --- sản xuất.} Về câu và đoạn:\cite{williams2016} là quyển cốt lõi,\cite{gopenswan1990} là bài báo mười trang đáng đọc kỹ nhất trong cả danh sách này. Về văn phong nói chung:\cite{zinsser2006} và\cite{pinker2014}; bản cực ngắn\cite{strunkwhite1999} nên đọc như một tài liệu lịch sử chứ không như luật. Về viết học thuật theo thể loại:\cite{swalesfeak2012} và\cite{booth2016}; về phê bình văn phong học thuật nặng nề,\cite{sword2012}. Về trình bày kỹ thuật và hình vẽ:\cite{alley2018}. Về khuôn câu để bắt đầu:\cite{graff2021} và\cite{manchesterphrasebank}. Về quy trình duy trì việc viết:\cite{silvia2007}.

\emph{Nhánh 4 --- lưu loát.} Không cần tài nguyên riêng; dùng lại tài liệu nhánh 1 với đồng hồ bấm giờ, và dùng chính bài báo bạn vừa đọc để nói lại trong ba phút.

\emph{Đọc và hiểu sâu (xuyên nhánh).} Về ba lượt đọc một bài báo:\cite{keshav2007}. Về đọc sách ở nhiều tầng:\cite{adler1972}. Về cấu trúc lập luận:\cite{toulmin1958}. Về siêu diễn ngôn và mức cam kết:\cite{hyland2005,hyland1998}. Về hàm ý và những gì người nói ngụ ý mà không nói:\cite{grice1975}. Về ẩn dụ:\cite{lakoff1980}. Về kết hợp từ và cụm:\cite{sinclair1991,lewis1993,wray2002,hoey2005}. Về khác biệt giữa các thể loại văn bản, đo bằng ngữ liệu:\cite{biber1999}.

\emph{Ghi chú và trí nhớ.} Về hệ ghi chú liên kết:\cite{ahrens2017}. Về cơ sở của việc ôn giãn cách và truy xuất:\cite{cepeda2006,roediger2006}.

\begin{cannho}
Đừng đọc danh mục này từ trên xuống. Cách dùng đúng: chọn \emph{một} quyển cho nhánh đang yếu nhất của bạn, đọc nó trong lúc chạy lộ trình, và chỉ lấy quyển thứ hai khi đã áp dụng xong quyển thứ nhất. Với phần lớn người ở trình độ của bạn, hai mục đáng bắt đầu là\cite{gopenswan1990} --- mười trang --- và\cite{williams2016}.
\end{cannho}

\section{Gốc rễ: vì sao lộ trình thất bại}

Các lộ trình tự học thất bại theo những cách khá giống nhau, và ba nguyên nhân dưới đây giải thích phần lớn.

\emph{Nguyên nhân một --- đặt cường độ theo tuần tốt nhất.} Người ta lập kế hoạch trong một buổi tối có động lực cao và ước lượng thời gian theo trạng thái đó. Xu hướng ước lượng lạc quan về thời gian hoàn thành một việc là một thiên lệch được mô tả rộng rãi, và nó không biến mất khi bạn biết về nó. Cách phòng duy nhất đáng tin là thiết kế cho ngày tệ nhất --- cột cuối của bảng ở mục 1.

\emph{Nguyên nhân hai --- không có phản hồi nên không thấy tiến bộ.} Ngôn ngữ tiến bộ chậm và không đều, nên cảm giác chủ quan gần như không dùng được làm tín hiệu. Không có chỉ số đo được, người học kết luận sai rằng mình không tiến và bỏ. Đây chính là lý do mục 3 tồn tại, và vì sao chỉ số thứ ba --- phản ứng nhanh nhất --- lại quan trọng hơn vẻ ngoài của nó.

\emph{Nguyên nhân ba --- mục tiêu không gắn với việc thật.} Một lộ trình tách rời khỏi công việc cạnh tranh trực tiếp với công việc về thời gian, và nó luôn thua. Một lộ trình gắn vào công việc thì không cạnh tranh --- đó là toàn bộ ý của \ptr{E/19} mục 5, và là lý do bảng ở mục 1 chỉ định các hoạt động mà bạn vốn đã làm.

Có một nguyên nhân thứ tư ít được nói tới và đáng nêu vì nó chạm vào chính quyển sách này: \emph{nhầm việc đọc về phương pháp với việc luyện tập}. Đọc một chương về cách viết câu rõ cho cảm giác tiến bộ ngay lập tức, còn viết mười lăm phút thì không. Cảm giác đó lệch, và nó khiến người ta đọc thêm sách thay vì làm. Nếu bạn nhớ một câu từ chương này, hãy nhớ câu đó: quyển sách này chỉ có giá trị bằng số lần bạn dùng nó trên một bài viết thật.

\begin{gocre}
\gr{Ước lượng thời gian lạc quan}{Vì sao kế hoạch luôn quá tay?}{Người ta lập kế hoạch ở trạng thái động lực cao \(\Rightarrow\) phải thiết kế cho ngày tệ nhất.}
\gr{Ôn giãn cách}{Ôn dồn hay ôn rải?}{Rải hiệu quả hơn\cite{cepeda2006} \(\Rightarrow\) lịch ôn ở mục 5.}
\gr{Hiệu ứng kiểm tra}{Đọc lại hay tự nhớ lại?}{Truy xuất củng cố mạnh hơn\cite{roediger2006} \(\Rightarrow\) ôn bằng phần Ôn nhanh, không bằng thân chương.}
\gr{Thiếu phản hồi}{Vì sao người học bỏ dù đang tiến?}{Ngôn ngữ tiến chậm và không đều; không có chỉ số thì cảm giác dẫn tới kết luận sai.}
\gr{Cạnh tranh thời gian}{Vì sao lộ trình tách rời công việc luôn thua?}{Nó cạnh tranh trực tiếp với công việc \(\Rightarrow\) gắn vào việc sẵn có (\ptr{E/19} mục 5).}
\gr{Đọc về phương pháp}{Vì sao đọc nhiều mà không khá?}{Đọc cho cảm giác tiến bộ tức thì mà không luyện gì \(\Rightarrow\) đo bằng sản phẩm, không bằng số chương đã đọc.}
\end{gocre}

\section{Liên hệ ngang}

\begin{lienhe}
\lh{Thể thao}{Chu kỳ tập luyện}{Chia giai đoạn có trọng tâm và có điểm kiểm tra --- cùng cấu trúc với ba giai đoạn ở mục 2, và cùng lý do: không thể tăng mọi mặt cùng lúc.}
\lh{Quản lý dự án}{Sản phẩm bàn giao theo mốc}{Sản phẩm cuối mỗi giai đoạn biến việc học thành việc có điểm kết thúc quan sát được.}
\lh{Kỹ thuật phần mềm}{Tích hợp liên tục}{Chạy đều mỗi ngày ở mức nhỏ tốt hơn dồn một đợt lớn --- cùng nguyên tắc với mức tối thiểu ở mục 1.}
\lh{Tài chính cá nhân}{Tiết kiệm đều đặn}{Số tiền nhỏ đều đặn thắng số tiền lớn thất thường; cùng toán học với việc không đứt chuỗi.}
\lh{Y tế dự phòng}{Tuân thủ điều trị}{Vấn đề trung tâm không phải phác đồ tốt nhất mà là phác đồ người ta theo được --- đúng lý do chọn năm giờ thay vì mười lăm.}
\lh{Nghiên cứu}{Sổ ghi thí nghiệm}{Ghi chỉ số hằng tuần cho bạn dữ liệu về chính mình, thay cho cảm giác --- cùng logic với việc ghi chép trong nghiên cứu.}
\end{lienhe}

\begin{traloi}
\tl{Năm giờ chia đều bốn nhánh: khoảng 75 phút đọc và nghe tài liệu bạn hiểu gần hết; 75 phút học có chủ đích (bảng thuật ngữ, ghi cụm, tra ngữ liệu); 90 phút sản xuất (nhật ký kỹ thuật 15 phút bốn lần, viết lại một đoạn); và 60 phút luyện lưu loát (đọc nhanh bấm giờ, nói lại một bài báo trong ba phút). Quan trọng không kém là cột mức tối thiểu cho ngày bận --- một trang không tra từ, ba mục thuật ngữ, năm câu viết, ba phút nói.}
\tl{Bằng chỉ số đo được, vì cảm giác tiến bộ trong ngôn ngữ không đáng tin theo cả hai chiều. Bốn chỉ số, mỗi cái ứng một nhánh: số lần khựng trên một trang (giảm; dưới hai là ngưỡng trôi chảy); số mục mới trong bảng thuật ngữ; số từ viết trong 15 phút không dừng --- chỉ số nhạy nhất, thường tăng rõ trong ba tới bốn tuần; và số cụm \emph{đã dùng} trong bài thật, chỉ số duy nhất phân biệt học thật với sưu tầm.}
\tl{Không làm bù. Làm bù biến việc lỡ thành một khoản nợ, và nợ tích lại làm người ta bỏ hẳn --- buổi đã lỡ thì mất, tiếp tục từ hôm nay. Trong tuần bận, hạ xuống mức tối thiểu chứ đừng bỏ hẳn, vì đứt chuỗi mới là thứ giết thói quen. Và nếu đã lỡ hơn hai tuần, đừng quay thẳng về cường độ cũ: chạy mức tối thiểu một tuần rồi mới lên lại, vì quay thẳng gần như luôn dẫn tới lần bỏ thứ hai.}
\tl{Giữ nguyên cấu trúc tuần nhưng gắn \emph{nội dung} vào một mục tiêu thật --- một bài báo đang viết, một buổi trình bày, một chương luận án. Lý do không phải động lực mà là chất lượng phản hồi: bài thật có người phản biện thật, và đó là thứ mọi bài luyện tự tạo đều thiếu. Đây cũng là lúc chuyển từ ``luyện tiếng Anh'' sang ``làm việc bằng tiếng Anh'', tức trạng thái đích của cả quyển sách.}
\tl{Đừng đọc lại tuần tự. Ba bước: ôn bằng phần \emph{Ôn nhanh} chứ không bằng thân chương --- che đáp án, tự trả lời trước, và lần đầu chỉ làm các mục mức 3; ôn theo lịch giãn --- sau một tuần, ba tuần, hai tháng, ghi ngày vào dòng ``Ôn gần nhất'' đầu mỗi tệp; và ôn theo triệu chứng --- khi gặp vấn đề thật, dùng bảng tra ngược ở đầu mỗi phần để nhảy thẳng tới chương liên quan. Cách thứ ba có giá trị nhất vì nó xảy ra đúng lúc bạn có nhu cầu.}
\end{traloi}

\begin{dongian}
Đến phần cuối cùng: làm gì vào thứ Hai tuần này.

Kế hoạch dưới đây tính cho \emph{năm giờ một tuần}. Con số này không phải tối ưu --- nó là con số mà một người đi làm giữ được trong nhiều tháng. Kế hoạch mười lăm giờ trông đẹp trên giấy và bị bỏ sau ba tuần, mà ba tuần thì không đủ để thay đổi gì.

Năm giờ chia làm bốn phần, tương ứng bốn nhánh:

Khoảng một tiếng rưỡi đọc những thứ bạn hiểu gần hết --- bài tổng quan, giáo trình trong ngành. Một tiếng rưỡi tra từ, ghi cụm, xây bảng thuật ngữ. Một tiếng rưỡi viết: mỗi ngày mười lăm phút, không dừng lại tra từ. Và một tiếng làm những việc dễ nhưng nhanh: đọc có bấm giờ, nói lại nội dung một bài báo trong ba phút.

Nhưng phần quan trọng nhất của kế hoạch không phải mấy con số đó. Nó là \emph{mức tối thiểu cho ngày bận nhất}: một trang đọc, ba mục thuật ngữ, năm câu viết, ba phút nói. Vì thứ làm người ta bỏ cuộc không phải làm ít --- mà là bỏ trắng một ngày, rồi ngày thứ hai.

Về chuyện lỡ nhịp, có một quy tắc đi ngược trực giác nhưng rất hiệu quả: \emph{đừng làm bù}. Làm bù biến buổi đã lỡ thành một món nợ, và nợ chồng lên thì người ta bỏ hẳn. Buổi đó mất rồi, tiếp tục từ hôm nay.

Cuối cùng, hãy đo bằng con số thay vì bằng cảm giác, vì cảm giác về tiến bộ ngôn ngữ rất hay sai. Bốn con số dễ đo: số lần bạn phải dừng lại trên một trang; số mục mới trong bảng thuật ngữ; số từ viết được trong mười lăm phút; và số cụm bạn đã thực sự \emph{dùng} trong một bài thật. Con số thứ ba nhúc nhích sớm nhất --- thường thấy rõ sau ba, bốn tuần.

Và một điều đáng nói về chính quyển sách này: đọc nó cho cảm giác tiến bộ ngay, còn viết mười lăm phút thì không cho cảm giác gì. Cảm giác đó lệch. Quyển sách này chỉ có giá trị bằng số lần bạn mở nó ra trong lúc đang sửa một bài viết thật.
\motcau{Thiết kế cho ngày tệ nhất, không phải ngày tốt nhất.}
\chuadongian{Mức tối thiểu đúng cho riêng bạn chỉ tìm được bằng cách thử vài tuần rồi điều chỉnh.}
\end{dongian}

\begin{onbai}
\ar[3]{Vì sao chọn năm giờ một tuần}{Đó là mức duy trì được nhiều tháng; mười lăm giờ bị bỏ sau ba tuần.}
\ar[3]{Tuần mẫu}{75 phút nhánh 1; 75 phút nhánh 2; 90 phút nhánh 3; 60 phút nhánh 4.}
\ar[3]{Mức tối thiểu ngày bận}{Một trang; ba mục thuật ngữ; năm câu; ba phút nói --- để chuỗi không đứt.}
\ar[3]{Ba giai đoạn}{Tuần 1--4 nền và đọc; 5--8 viết rõ; 9--12 sắc thái và tự động hoá.}
\ar[3]{Vì sao đọc đi trước viết}{Viết cần nguyên liệu và cần mẫu.}
\ar[3]{Bốn chỉ số}{Lần khựng mỗi trang; mục thuật ngữ mới; số từ trong 15 phút; số cụm \emph{đã dùng}.}
\ar[2]{Chỉ số nhạy nhất}{Số từ viết trong 15 phút --- tăng rõ trong 3--4 tuần.}
\ar[2]{Chỉ số phân biệt học thật với sưu tầm}{Số cụm đã dùng, không phải số cụm đã ghi.}
\ar[2]{Ba quy tắc khi lỡ nhịp}{Không làm bù; hạ xuống mức tối thiểu; lỡ hơn hai tuần thì vào lại bằng mức tối thiểu.}
\ar[2]{Bẫy chuyển sang nhánh dễ}{Khi mệt người ta dồn vào đọc và nghe --- tuần đó là tuần mất cân bằng, không phải tuần nhẹ.}
\ar[2]{Sau 12 tuần}{Giữ cấu trúc, gắn nội dung vào một mục tiêu thật để có phản hồi thật.}
\ar[2]{Ba cách ôn lại quyển này}{Ôn bằng phần Ôn nhanh; ôn theo lịch giãn; ôn theo triệu chứng bằng bảng tra ngược.}
\ar[1]{Bốn nguyên nhân lộ trình thất bại}{Đặt cường độ theo tuần tốt nhất; thiếu chỉ số; tách rời việc thật; và nhầm đọc về phương pháp với luyện tập.}
\end{onbai}

\begin{hoidap}
\qa{Tôi có ít hơn năm giờ một tuần. Cắt gì?}{Đừng cắt một nhánh; hãy thu nhỏ cả bốn. Với ba giờ: 45 phút mỗi nhánh. Với hai giờ: 30 phút mỗi nhánh. Lý do quan trọng --- cắt hẳn một nhánh phá đúng thứ mà \ptr{E/19} nói là cốt lõi, trong khi làm ít ở cả bốn vẫn giữ được cân bằng và chỉ làm chậm tiến độ. Nếu buộc phải ưu tiên, hãy giữ nhánh 3 nguyên vẹn và thu nhỏ ba nhánh kia, vì nhánh 3 là nhánh bạn có nhiều khả năng đang thiếu nhất.}
\qa{Nhật ký kỹ thuật thì viết cái gì?}{Bất cứ thứ gì thuộc công việc của bạn: hôm nay bạn thử gì, kết quả ra sao, bạn nghĩ vì sao nó không chạy, bước tiếp theo là gì. Ưu điểm lớn là bạn không phải nghĩ nội dung --- nội dung đã có sẵn trong đầu bạn, nên toàn bộ nỗ lực dồn vào việc diễn đạt. Và nó có ích ngoài tiếng Anh: sau ba tháng bạn có một bản ghi công việc mà bạn sẽ dùng lại khi viết bài.}
\qa{Có nên tìm người sửa bài cho tôi không?}{Rất nên nếu tìm được, vì phản hồi là mắt xích yếu nhất của người tự học (\ptr{E/19} mục 4). Nhưng đừng chờ có người mới bắt đầu. Ba nguồn phản hồi thay thế chạy được ngay: đối chiếu bài của bạn với một bài mẫu cùng thể loại; tra bằng ngữ liệu ngành; và hỏi mô hình ngôn ngữ \emph{chỗ nào} khó hiểu theo ranh giới ở \ptr{E/20}. Khi tìm được người đọc, hãy hỏi họ ``chỗ nào bạn phải đọc lại'' chứ không hỏi ``bài có hay không''.}
\qa{Lộ trình này có hợp với người muốn thi chứng chỉ không?}{Chỉ một phần. Các kỳ thi chuẩn hoá đo một tập kỹ năng hẹp hơn và có định dạng riêng, nên chúng đòi luyện đề --- việc mà lộ trình này không có. Phần nền thì vẫn dùng được, đặc biệt là nhánh 3 và nhánh 4. Nếu bạn cần cả hai, hãy chạy lộ trình này làm nền và thêm bốn tới sáu tuần luyện định dạng ngay trước kỳ thi; đừng luyện đề suốt nhiều tháng, vì nó cải thiện điểm mà cải thiện rất ít khả năng thật.}
\qa{Tôi nên đọc lại quyển này bao lâu một lần?}{Đừng nghĩ theo lịch mà nghĩ theo triệu chứng --- bước ba ở mục 5. Ngoại lệ là hai chương nên ôn theo lịch vì bạn dùng chúng liên tục mà lại dễ quên chi tiết: \ptr{C/12} và \ptr{C/14}. Với chúng, ôn phần Ôn nhanh mỗi tháng một lần trong nửa năm đầu. Các chương còn lại để dành cho lúc gặp đúng vấn đề mà chúng xử lý.}
\qa{Làm sao biết mình đã ``đủ giỏi'' và có thể dừng?}{Câu hỏi này có một câu trả lời thực dụng: khi tiếng Anh không còn là thứ giới hạn công việc của bạn. Cụ thể hơn --- bạn đọc tài liệu ngành mà không khựng quá hai lần một trang; bạn viết một phần bài mà không phải dừng lại tra cấu trúc; và phản hồi bạn nhận về bài viết nói về \emph{nội dung} chứ không về ngôn ngữ. Khi đạt trạng thái đó, việc luyện tập tách riêng có thể dừng --- cái tiếp tục là việc dùng, và nó tự duy trì.}
\qa{Trong danh mục tài nguyên, nếu chỉ đọc hai thứ thì đọc gì?}{Bài báo mười trang của Gopen và Swan\cite{gopenswan1990} và quyển của Williams\cite{williams2016}. Lý do: chúng phủ đúng hai nguyên tắc cho phần lớn mức cải thiện trong viết --- chủ thể--hành động và cũ trước mới sau --- và chúng dạy bằng ví dụ sửa cụ thể chứ không bằng lời khuyên chung. Nếu bạn muốn thêm một quyển thứ ba cho việc học từ vựng có hệ thống, lấy\cite{webbnation2017}.}
\end{hoidap}

\begin{baitap}
\bt[1]{Viết ra tuần mẫu của riêng bạn, kèm cột mức tối thiểu}{Tuần này}{Gắn mỗi nhánh vào một mốc sẵn có trong ngày.}
\bt[1]{Đo bốn chỉ số lần đầu và ghi lại}{Hôm nay}{Đây là mốc gốc; không có nó thì không so được.}
\bt[2]{Chạy trọn giai đoạn 1 bốn tuần và nộp sản phẩm cho chính mình}{Tuần 1--4}{Bảng thuật ngữ 60 mục và một bản đồ lập luận.}
\bt[2]{Đo lại bốn chỉ số mỗi tuần trong mười hai tuần}{12 tuần}{Vẽ ra một đường; chỉ số 3 sẽ nhúc nhích trước.}
\bt[2]{Viết quy tắc lỡ nhịp của riêng bạn ra giấy trước khi bắt đầu}{Hôm nay}{Quyết định trước lúc bình tĩnh, không lúc đang mệt.}
\bt[3]{Hoàn thành sản phẩm của cả ba giai đoạn}{12 tuần}{Sản phẩm cuối: một bài viết hoàn chỉnh đã sửa bốn lượt.}
\bt[3]{Sau 12 tuần, chọn một mục tiêu thật và thiết kế lại nội dung bốn nhánh quanh nó}{Tuần 13}{Cấu trúc giữ nguyên, nội dung gắn vào bài thật.}
\end{baitap}

\section*{Kết chương và đọc tiếp}

Quyển sách khép lại ở đây với một lịch cụ thể, bốn chỉ số, và một danh mục tài nguyên. Nhưng phần quan trọng nhất của chương này là một nhận xét về chính nó: đọc về phương pháp cho cảm giác tiến bộ ngay lập tức, còn viết mười lăm phút thì không --- và cảm giác đó lệch.

Nếu bạn chỉ làm một việc sau khi gấp quyển này lại, hãy làm việc nhỏ nhất trong bảng ở mục 1: viết năm câu tiếng Anh về công việc hôm nay, không dừng lại tra từ. Nó nhỏ tới mức không có lý do gì để không làm, và nó thuộc đúng nhánh mà bạn nhiều khả năng đang bỏ trống.

Phần \ptr{Z/97} sau đây là danh sách đọc thêm có chú giải, \ptr{Z/98} là bảng thuật ngữ của cả quyển, và \ptr{Z/99} là vài lời kết.
\ifSubfilesClassLoaded{\printbibliography[title={Nguồn}]}{}
\end{document}
